\section{Conclusion}
\label{sec:conclusion}

We proposed \(\protocol\), a new protocol for verifying matrix multiplication.
The key idea is to move the expensive vector-matrix products outside the SNARK
circuit and replace them with sampling-based proximity tests over linear
codewords.  The prover first encodes and commits to the matrices, then computes
the Freivalds intermediate vectors using a verifier-sampled random challenge.
The prover then encodes and commits to these intermediate vectors, and the
SNARK circuit proves only local fold relations at randomly sampled positions.
Merkle openings and CP-Link proofs ensure that the values used inside the
CP-SNARK are consistent with the committed data fixed before the sampling
challenges are known.

The security analysis follows this structure.  Assuming that \(\rt_{ABC}\)
binds to valid row-wise encodings of \(\mat{A},\mat{B},\mat{C}\), any prover that
convinces the verifier of a false statement \(\mat{A}\mat{B}\neq\mat{C}\) must
either hit the Freivalds failure event, avoid every sampled codeword
disagreement, or break one of the backend assumptions of \(\protocol\).  These
assumptions include CP-SNARK soundness, commitment binding, CP-Link consistency,
and code-consistency soundness.

Our implementation confirms the intended tradeoff.  For a fixed number of
queries, the main arithmetic relation has \(\mathcal{O}(tk)\) constraints
rather than the \(\mathcal{O}(k^2)\) constraints of a Freivalds-in-SNARK
baseline.  In the matrix benchmark at \(k=4{,}096\), \(\protocol\) reduces the
constraint count by \(30.3\times\) and proof generation time by \(8.34\times\),
while keeping verification below \(10\) ms.  The GPT-2 workload shows that
\(\protocol\) can also be applied to matrix-heavy verifiable AI computation.
When many multiplication claims appear together, the gains are smaller but
remain meaningful.

The main remaining cost of \(\protocol\) is proof size.  Since sampled Merkle
openings are included in the proof, this overhead becomes more visible as the
number of matrix-multiplication claims increases.  A natural next step is to
optimize the opening layer for workloads in which many matrix products appear
in sequence.  In particular, batching or sharing Merkle openings and CP-Link
proofs across multiple products may reduce proof size and verification overhead
while preserving the linear in-circuit cost that \(\protocol\) provides.
